
\documentclass[11pt]{article}
\usepackage[margin=1in]{geometry}
\usepackage{hyperref}
\usepackage{enumitem}
\usepackage{titlesec}
\usepackage{graphicx}
\usepackage{setspace}

\title{ExamProfile AI \\ \large Professor-Modeled Exam Generation and Weakness Analytics for STEM Courses}
\author{Project Specification}
\date{\today}

\begin{document}
\maketitle
\tableofcontents
\newpage

\section{Overview}
ExamProfile AI is a STEM-focused, web-first application for exam simulation and weakness analytics. It ingests course materials such as lecture slides, homework sets, and prior exams to generate course-specific practice problem sheets, simulated exams, and answer keys. The primary delivery surface is intended to be a web application, while keeping product language flexible enough to support a desktop client in the future.

Unlike generic AI study tools, ExamProfile AI builds a structured \textbf{Professor Profile}: an evidence-weighted model of instructor tendencies and course assessment patterns (e.g., recurring topics, typical question formats, exam structure, and estimated difficulty tendencies). Topic emphasis is treated as informative rather than deterministic; if course evidence strongly emphasizes a topic (such as probability), generated exams should generally reflect that tendency while preserving realistic variety. Stronger evidence signals (e.g., past exams, practice tests, explicit instructor guidance, or user-provided prompts/comments) are intended to be weighted more heavily than weak signals (e.g., a single homework problem).

Generation and evaluation are intended to be multi-stage and quality-controlled (generation, review, validation, and refinement), rather than a one-shot ``generate an exam'' prompt. Students submit answers and receive diagnostically useful feedback (not only correct/incorrect), including structured error classifications that support concept-level weakness analytics and targeted follow-up practice.

\section{Problem}
Students preparing for exams face several recurring issues:

\begin{itemize}
\item Generic AI-generated practice is rarely course-specific; it does not reliably reflect an instructor's exam structure, pacing, grading style, or concept emphasis as inferred from real course materials.
\item Typical LLM study tools often produce shallow variants that restate known problems with minor wording or numeric changes, which does not adequately prepare students for new but conceptually aligned exam questions.
\item Existing tools frequently fail to reproduce realistic exam construction (e.g., section structure, question-type distribution, and difficulty progression) and may generate multiple-choice questions with repetitive or predictable answer-position patterns.
\item Feedback is commonly limited to correctness labels or generic explanations and is not diagnostically structured (e.g., it does not identify wrong method selection versus formula misuse versus computation versus interpretation errors).
\end{itemize}

Because of this, students often study inefficiently and cannot reliably simulate real exam conditions.

\section{Product Vision}
The goal of ExamProfile AI is to create a system that:
\begin{itemize}
\item Ingests real course materials to learn course scope and assessment patterns
\item Builds and maintains a structured, evidence-weighted Professor Profile that informs (but does not rigidly dictate) generation
\item Produces simulated exams and practice sets that are intended to be novel-but-aligned with instructor style, structure, and concept emphasis (avoiding obvious paraphrases unless explicitly requested)
\item Grades submissions with diagnostically useful, structured feedback beyond correct/incorrect, identifying what was correct, which method was misapplied, and what conceptual weakness is suggested
\item Iterates a preparation loop: ingest materials $\rightarrow$ model tendencies $\rightarrow$ generate practice $\rightarrow$ diagnose performance $\rightarrow$ update weakness signals $\rightarrow$ regenerate targeted practice
\item Uses a multi-stage generation, review, validation, and refinement pipeline to improve realism and reduce brittle one-shot outputs (later sections specify these mechanisms in detail)
\end{itemize}

The system is designed to operate as an \textbf{exam simulation engine} grounded in course evidence, not as a generic chatbot tutor and not as a replacement for instruction.

\section{Target Users}
Primary users include:
\begin{itemize}
\item College students in quantitative, structured STEM courses preparing for midterms and finals
\item Students who can provide course materials (slides, homework, prior exams) and want course-specific, professor-aligned exam simulation
\item Students who need diagnostically structured feedback on methods, formulas, and reasoning patterns (not only correctness labels)
\end{itemize}

Initial focus domains include:
\begin{itemize}
\item Statistics and probability
\item Applied mathematics
\item Chemistry and physics
\item Computational biology
\item Computer science courses with structured assessments
\end{itemize}

The initial focus is optimized for structured, quantifiable subjects where methods, formulas, and reasoning patterns can be evaluated consistently; it is not initially targeted at essay-heavy or primarily subjective disciplines.

\section{Key Differentiators}
\begin{enumerate}
\item \textbf{Professor Modeling}
Builds a structured Professor Profile capturing instructor tendencies and course patterns (topic emphasis, question-type distribution, exam structure, and estimated difficulty tendencies), using an evidence-weighted approach that prioritizes stronger signals (e.g., past exams) when available.

\item \textbf{Distribution-Aware Exam Generation}
Generates simulated exams informed by observed topic emphasis and structure without mechanically matching exact topic percentages. The system aims to create new, non-trivially distinct questions that are aligned with instructor style and avoids shallow paraphrases by default. For multiple-choice questions, generation aims for intentionally balanced, realistic construction (including avoiding repetitive or predictable answer-position patterns).

\item \textbf{Multi-Stage Quality-Controlled Generation}
Uses a layered generation pipeline rather than a single LLM call. Draft questions and exams can be reviewed, validated, and refined before delivery to improve reliability and to reduce shallow reformulations and recurring formatting artifacts.

\item \textbf{Scope-Aware Filtering}
Supports constraints such as ``Only review Notes 3--8'' or ``Focus on hypothesis testing'' to prevent out-of-scope generation and to focus practice on a defined portion of the course.

\item \textbf{Concept-Level Weakness Analytics}
Provides diagnostically structured grading signals (e.g., wrong method selection, formula misuse, computational mistakes, interpretation errors) to identify concept-level weaknesses. Weakness analytics guide targeted regeneration of follow-up practice, reinforcing the exam-preparation loop rather than producing one-off problem sets.
\end{enumerate}

\section{Core User Flow}

\subsection{Step 1: Upload Materials}
Users upload:
\begin{itemize}
\item Lecture slides
\item Homework sets
\item Prior exams (optional)
\end{itemize}

Users may also specify scope constraints such as topic ranges.

\subsection{Step 2: Build Professor Profile}
The system:
\begin{itemize}
\item Parses documents
\item Identifies topics and question types
\item Builds a Professor Profile
\end{itemize}

\subsection{Step 3: Insights}
The system displays:
\begin{itemize}
\item extracted topics
\item topic emphasis
\item inferred exam structure
\end{itemize}

\subsection{Step 4: Practice Materials}
The system generates:
\begin{itemize}
\item Practice problem sheets
\item Answer keys with explanations
\end{itemize}

\subsection{Step 5: Practice Exam Generation}
Users select:
\begin{itemize}
\item MCQ only
\item FRQ only
\item Mixed mode with custom ratio
\end{itemize}

The system generates a simulated exam.

\subsection{Step 6: Submission and Grading}
Students submit answers through the interface. 
The system grades responses and classifies errors.

\subsection{Step 7: Weakness Analytics}
The system produces:
\begin{itemize}
\item concept mastery heatmaps
\item error classifications
\item targeted practice recommendations
\end{itemize}

\section{Document Parsing Layer}
Traditional PDF parsers struggle with scientific documents containing equations and multi-column layouts.

The system uses layout-aware tools:

\begin{itemize}
\item \textbf{Docling (IBM)} --- Converts complex PDFs into structured Markdown suitable for LLM pipelines.
\item \textbf{Marker} --- Alternative parser with strong support for mathematical documents.
\item \textbf{Pandoc} --- Used for converting generated Markdown into final PDF outputs.
\end{itemize}

\section{AI and Retrieval Pipeline}

\subsection{Document Parsing}
Uploaded PDFs are converted into structured Markdown using Docling.

\subsection{Chunking}
Documents are split into semantic chunks representing sections or examples.

\subsection{Embeddings}
Each chunk is converted into an embedding vector representing semantic meaning.

\subsection{Retrieval}
Relevant chunks are retrieved based on user constraints or exam generation queries.

\subsection{Professor Profile Construction}
Topic frequencies and question patterns are extracted from retrieved materials.

\subsection{Exam Generation}
The LLM generates exam questions based on:
\begin{itemize}
\item retrieved course materials
\item the professor profile
\item exam format settings
\end{itemize}

\subsection{Structured Grading}
Student responses are evaluated using rubric-based grading logic.

\subsection{Adaptive Regeneration}
Weak concepts are fed back into the generation pipeline to produce targeted follow-up questions.

\section{Weakness Dashboard}
The dashboard provides:

\begin{itemize}
\item Concept mastery heatmaps
\item Error type classification
\item Performance trends
\item Targeted practice recommendations
\end{itemize}

Example error categories include:

\begin{itemize}
\item Wrong method selection
\item Formula misuse
\item Computational mistakes
\item Interpretation errors
\end{itemize}

\section{Technology Stack}

\subsection{Frontend}
\begin{itemize}
\item React
\item TypeScript
\item Tailwind CSS
\item Shadcn/ui
\item Recharts for analytics dashboards
\end{itemize}

\subsection{Backend}
\begin{itemize}
\item Python
\item FastAPI
\end{itemize}

\subsection{Document Parsing}
\begin{itemize}
\item Docling
\item Marker
\item Pandoc
\end{itemize}

\subsection{AI and Retrieval}
\begin{itemize}
\item Gemini 1.5 Flash API
\item LlamaIndex for orchestration
\item Retrieval-Augmented Generation (RAG)
\end{itemize}

\subsection{Vector Database}
\begin{itemize}
\item ChromaDB (local development)
\item Optional Pinecone for managed cloud deployment
\end{itemize}

\subsection{Primary Database}
\begin{itemize}
\item Supabase (PostgreSQL + pgvector)
\end{itemize}

\subsection{Deployment}
\begin{itemize}
\item Frontend: Vercel
\item Backend: Render or Railway
\item Database: Supabase
\end{itemize}

\section{Risks and Constraints}

\begin{itemize}
\item Parsing quality of complex PDFs
\item Reliability of FRQ grading
\item Accurate mapping of topic scope constraints
\item Preventing scope creep beyond MVP
\end{itemize}

\section{One-Sentence Pitch}
ExamProfile AI is a web-first exam simulation and weakness analytics platform that learns instructor tendencies from real course materials, generates novel-but-aligned simulated exams, and provides diagnostically structured feedback to drive targeted preparation.

\end{document}
